% Auto-generated by make_appendix_tables.py — do not hand-edit.
\begin{table}[t]
\caption{Shipborne case inventory and residuals: the four collocated overpasses of the two cruises.}
\label{tab:ship-cases}
\footnotesize
\begin{tabular}{llrrrrrrr}
\tophline
Date & Cruise & $n$ & $\tilde{d}_{\mathrm{cld}}$ & $X_{\mathrm{CO_2}}^{\mathrm{ship}}$ & $X_{\mathrm{CO_2}}^{\mathrm{B11}}$ & $X_{\mathrm{CO_2}}^{\mathrm{DE}}$ & res$_{\mathrm{B11}}$ & res$_{\mathrm{DE}}$ \\
\middlehline
2019-06-09 & R/V Sonne SO268 & 586 & 2.2 & 412.35 & 413.44 & 413.52 & +1.09 & +1.17 \\
2019-06-14 & R/V Sonne SO268 & 243 & 4.9 & 411.91 & 412.13 & 412.35 & +0.22 & +0.44 \\
2019-06-22 & R/V Sonne SO268 & 644 & 41.3 & 411.22 & 411.98 & 412.01 & +0.77 & +0.80 \\
2021-03-15 & R/V Mirai MR21-01 & 21 & 2.9 & 415.25 & 417.13 & 417.55 & +1.89 & +2.30 \\
\bottomhline
\end{tabular}
\\[3pt]\parbox{\linewidth}{\footnotesize Residuals: OCO-2 minus shipborne EM27/SUN column (ppm). $\tilde{d}_{\mathrm{cld}}$: median nearest-cloud distance (km). The shipborne columns carry an absolute-scale limitation (Sect.~4.6); the clear-day case 2019-06-22 serves as a negative control (Fig.~\ref{app-fig:ship_far_cld}).}
\end{table}
